
\begin{frame}{Historical and Representation Bias}

    \begin{columns}[T]
        \begin{column}{0.48\textwidth}
            \textbf{Historical bias}
            \vspace{0.3em}

            Present even with perfect sampling and perfect measurement, because the world the data
            records is itself unequal.

            \vspace{0.6em}
            \begin{itemize}\small
                \setlength\itemsep{0.35em}
                \item An image search for ``CEO'' that reflects who currently holds the job.
                \item A loan model trained on approvals from an era of discriminatory lending.
            \end{itemize}

            \vspace{0.6em}
            \begin{tcolorbox}[colback=gray!8, boxrule=0pt]
                \small More data does not fix this. It is a question about the objective, not
                the sample.
            \end{tcolorbox}
        \end{column}
        \begin{column}{0.48\textwidth}
            \textbf{Representation bias}
            \vspace{0.3em}

            The development sample under-represents part of the use population, so the model
            generalises poorly there.

            \vspace{0.6em}
            \begin{itemize}\small
                \setlength\itemsep{0.35em}
                \item A dermatology model trained mostly on light skin.
                \item A speech model trained on one accent.
                \item Gender Shades: existing face benchmarks were 79.6\% and 86.2\%
                      lighter-skinned subjects.
            \end{itemize}

            \vspace{0.6em}
            \begin{tcolorbox}[colback=gray!8, boxrule=0pt]
                \small More \emph{targeted} data does help --- but you must first know the
                composition of your data.
            \end{tcolorbox}
        \end{column}
    \end{columns}
\end{frame}
